% =====================================================================
%  TABLA — atajos para escribir tablas sin \begin...\end
%
%  \tabla{contenido}                -> tabla normal (tablaminimal)
%  \tabla[lcc]{contenido}           -> con las columnas fijadas
%  \tablacolor{contenido}           -> tabla con cabecera de color
%  \tablarejilla{contenido}         -> tabla con rejilla completa
%  \ficha{contenido}                -> ficha de clave/valor
%
%  Las columnas se detectan solas (contando las "&" de las filas), asi
%  que normalmente no hay que escribir nada mas:
%
%     \tabla{
%       \tablacab{Característica} & \tablacab{VPS} & \tablacab{IaaS} \\
%       \midrule
%       Contratación & Cuota fija & Pago por uso \\
%     }
% =====================================================================

\ExplSyntaxOn

\tl_new:N \l__tabla_spec_tl
\int_new:N \l__tabla_int
\int_new:N \l__tabla_amp_int
\int_new:N \l__tabla_filas_int
\newcount\tabladivisor
\seq_new:N \l__tabla_amp_seq
\seq_new:N \l__tabla_filas_seq

% --- cuenta los items de una secuencia ---
\cs_new_protected:Npn \tabla_contar:nn #1 #2
  {
    \int_zero:N \l__tabla_int
    \seq_map_inline:Nn #1 { \int_incr:N \l__tabla_int }
    \int_set:Nn #2 { \l__tabla_int }
  }

% --- nº de columnas (deja el especifico en \l__tabla_spec_tl) ---
\cs_new_protected:Npn \tabla_columnas:n #1
  {
    % las celdas se separan por "&" y las filas por "\\" (el mismo token
    % en modo normal y en modo expl3)
    \seq_set_split:Nnn \l__tabla_amp_seq { & } { #1 }
    \seq_set_split:Nnn \l__tabla_filas_seq { \\ } { #1 }
    \tabla_contar:nn \l__tabla_amp_seq \l__tabla_amp_int
    \tabla_contar:nn \l__tabla_filas_seq \l__tabla_filas_int
    % las "&" de menos (la primera celda no lleva "&")
    \int_sub:Nn \l__tabla_amp_int { 1 }
    % las filas de menos (la ultima puede no cerrar con "\\")
    \int_sub:Nn \l__tabla_filas_int { 1 }
    % columnas = "&" por fila + 1  (con \divide, que es de TeX puro)
    \tabladivisor=\int_use:N \l__tabla_amp_int
    \divide\tabladivisor by \int_use:N \l__tabla_filas_int
    \advance\tabladivisor by 1
    \ifnum\tabladivisor<1 \tabladivisor=1 \fi
    % primera columna a la izquierda, las demas centradas
    \tl_clear:N \l__tabla_spec_tl
    \tl_put_right:Nn \l__tabla_spec_tl { l }
    \loop
      \tl_put_right:Nn \l__tabla_spec_tl { c }
      \advance\tabladivisor by -1
    \ifnum\tabladivisor>1
    \repeat
  }

% --- \tablaimpl{columnas}{contenido} ---
\cs_new_protected:Npn \tablaimpl #1 #2
  {
    \group_begin:
      \tl_set:Nn \l__tabla_spec_tl { #1 }
      \tl_if_blank:eTF { \l__tabla_spec_tl }
        { \tabla_columnas:n {#2} }
        { }
      \edef\tablaenv{ \noexpand\begin{tablaminimal}{ \l__tabla_spec_tl } }
      \tablaenv
        #2%
      \end{tablaminimal}
    \group_end:
  }

% --- \tablacolorimpl{columnas}{contenido} ---
\cs_new_protected:Npn \tablacolorimpl #1 #2
  {
    \group_begin:
      \tl_set:Nn \l__tabla_spec_tl { #1 }
      \tl_if_blank:eTF { \l__tabla_spec_tl }
        { \tabla_columnas:n {#2} }
        { }
      \edef\tablaenv{ \noexpand\begin{tablacabecera}{ \l__tabla_spec_tl } }
      \tablaenv
        #2%
      \end{tablacabecera}
    \group_end:
  }

% --- \tablarejillaimpl{columnas}{contenido} ---
\cs_new_protected:Npn \tablarejillaimpl #1 #2
  {
    \group_begin:
      \tl_set:Nn \l__tabla_spec_tl { #1 }
      \tl_if_blank:eTF { \l__tabla_spec_tl }
        { \tabla_columnas:n {#2} }
        { }
      \edef\tablaenv{ \noexpand\begin{tablacuadricula}{ | \l__tabla_spec_tl | } }
      \tablaenv
        #2%
      \end{tablacuadricula}
    \group_end:
  }

% --- \fichaimpl{ancho}{contenido} ---
\cs_new_protected:Npn \fichaimpl #1 #2
  {
    \group_begin:
      \tl_set:Nn \l__tabla_spec_tl { #1 }
      \tl_if_blank:eTF { \l__tabla_spec_tl }
        { \tl_set:Nn \l__tabla_spec_tl { 4cm } }
        { }
      \edef\tablaenv{ \noexpand\begin{fichaclave}{ \l__tabla_spec_tl } }
      \tablaenv
        #2%
      \end{fichaclave}
    \group_end:
  }

% alias sin dos puntos, para poder llamarlos desde el modo normal
\cs_new_eq:NN \tablaImpl        \tablaimpl
\cs_new_eq:NN \tablacolorImpl   \tablacolorimpl
\cs_new_eq:NN \tablarejillaImpl \tablarejillaimpl
\cs_new_eq:NN \fichaImpl        \fichaimpl

\ExplSyntaxOff

% --- los comandos que se escriben (con argumento opcional) ---
\newcommand{\tabla}[2][]{%
  \begingroup
  \tablaImpl{#1}{#2}%
  \endgroup}

\newcommand{\tablacolor}[2][]{%
  \begingroup
  \tablacolorImpl{#1}{#2}%
  \endgroup}

\newcommand{\tablarejilla}[2][]{%
  \begingroup
  \tablarejillaImpl{#1}{#2}%
  \endgroup}

\newcommand{\ficha}[2][]{%
  \begingroup
  \fichaImpl{#1}{#2}%
  \endgroup}
